\documentclass[10pt,a4paper]{amsart}
\usepackage[margin=19mm]{geometry}
\usepackage{amsmath,amssymb,mathtools}
\usepackage[scr=rsfs,cal=euler]{mathalfa}
\usepackage{xfrac}
\usepackage[unicode,hidelinks,bookmarks=false]{hyperref}
\newcommand{\ie}{i.e.\ }
\newcommand{\cf}{cf.\ }
\newcommand{\resp}{respectively\ }
\newcommand{\Subsection}{\subsection}
\providecommand{\nobreakdash}{\nobreak\hbox{-}}
\setlength{\emergencystretch}{2em}
\multlinegap0pt
\usepackage{bm}
\usepackage{bbm}
\usepackage{dsfont}
\usepackage{xspace}
\usepackage{cancel}
\usepackage{mathtools}
\usepackage{amsthm}
\makeatletter
\@ifundefined{prop}{\newtheorem{prop}{Proposition}}{}
\makeatother
\makeatletter
\expandafter\ifx\csname dashlist\endcsname\relax
\newenvironment{dashlist}{\begin{list}{--}{}}{\end{list}}
\fi
\expandafter\ifx\csname plainlist\endcsname\relax
\newenvironment{plainlist}{\begin{list}{}{}}{\end{list}}
\fi
\expandafter\ifx\csname tight\endcsname\relax
\newenvironment{tight}{\begin{list}{$\bullet$}{\setlength{\itemsep}{-2mm}}}{\end{list}}
\fi
\makeatother
\title[Computing a Connection Matrix and Persistence Efficiently fr]{Computing a Connection Matrix and Persistence Efficiently from a Morse Decomposition}
\author{Tamal K. Dey, Micha{\l} Lipi\'nski, Andrew Haas}
\date{Source: arXiv:2502.19369v2 (CC BY 4.0)}
\begin{document}
\maketitle

\noindent\emph{Source and licence.} Computing a Connection Matrix and Persistence Efficiently from a Morse Decomposition. Version arXiv:2502.19369v2 (CC BY 4.0), distributed under the Creative Commons Attribution 4.0 International licence, as recorded in the arXiv licence metadata for this exact version and shown on the abstract page https://arxiv.org/abs/2502.19369v2. Full rights notice: licenses/arxiv-2502.19369-CC-BY-4.0.txt.\par\smallskip

\noindent\textit{Editorial scope.} Counted unit: the Prop whose source label is \texttt{prop:morsefixed} in the article (source0.tex lines 1763--1770, source label \texttt{prop:morsefixed}), delivered together with the article's own complete original proof of it (source0.tex lines 1771--1807, one \texttt{proof} environment of 37 lines). No mathematics is written, repaired or re-derived: statement and proof are the article's own text line for line. Required context retained uncounted: prop:morsepath (lines 1735--1740). Every definition, display and label of those lines is retained unchanged. Omitted from this package and not redistributed: 1--1734, 1741--1762, 1808--2269. Nothing inside a retained excerpt is omitted or altered: every retained body is delivered exactly as the article's lines are written, so no omission receipt of any kind accompanies this package. Citations: the retained excerpts carry no citation command at all, so no bibliography entry of the article is transcribed and no reference is redistributed. Reference adaptation: none was needed and none was made; every reference inside the delivered unit resolves inside it, so no unresolved reference or citation is produced. Printed slips kept literal: no printed word, symbol, hypothesis or number of the counted statement or of its proof was altered. Layout and class: the article's own class is \texttt{article} with packages including geometry, xcolor, tikz-cd, tikz, adjustbox, ulem, multirow, amsmath, amsthm, thmtools, hyperref, float, xy, array; the wrapper is the lane's pinned \texttt{amsart} editorial wrapper at 10pt on A4, which restates the theorem-like environments the retained text needs and adds the article's own macro definitions that the retained text needs (none), with extra wrapper packages (bm, bbm, dsfont, xspace, cancel, mathtools). The delivered numbering is the unit's own. Assets: no retained span contains a figure environment, a graphics-inclusion command, a PDF-page inclusion, a table or a TikZ picture; no image file is copied into this package, the package declares no asset dependency and no image file is redistributed with it. Licence: version arXiv:2502.19369v2 of this article is distributed under the Creative Commons Attribution 4.0 International licence, as recorded in the arXiv licence metadata for this exact version and shown on the abstract page https://arxiv.org/abs/2502.19369v2. A full rights notice is delivered at licenses/arxiv-2502.19369-CC-BY-4.0.txt. Characters: every retained excerpt is pure ASCII, so the delivered engine, XeLaTeX with its default Latin Modern text font, reports no missing character.
\begin{prop}
    For every stage of {\sc ConMat}, $A_i[l,j]\neq 0$ implies $[l]_P\leq_P [j]_P$.
    Moreover, if {\sc ConMat} adds column $k$ to column $j$ during its execution, then
    $[k]_P\leq_P [j]_P$.
    \label{prop:morsepath}
\end{prop}

\begin{prop}
Let $A_{out}$ and $A'_{out}$ be the matrices produced by {\sc ConMat} after
reducing Morse fixed input matrices $A$ and $A'$ respectively. Then,
  $A_{out}[i,j]=A'_{out}[i',j']$ where $i=_\sigma i'$ and $j=_\sigma j'$. 
  Furthermore, if $\ell$ is the pivot row of a homogeneous column $j$ in $A_{out}$,
  then $\ell'=_\sigma \ell$ is the pivot row of the homogeneous column $j'$ in $A_{out}'$.
\label{prop:morsefixed}
\end{prop}

\begin{proof}
    We prove the claim by inducting on $j$ where $A_j$ represents
    the processed matrices $A$ up to
    column $j$.
    
    The first column ($j=0$) corresponds to a vertex,
    thus has no non-zero entry in both matrices $A_{out}$ and $A'_{out}$.
    % \michal{The above sentence would be true for any Lefschetz complex, in particular, the first column would not have to correspond to a vertex. 
    %     This made me think, is there any place in the paper where we make use of the assumption (made at the beginning of Section \ref{sec:Cvec}) that $K$ is a simplicial complex?
    %     Maybe there is no point of changing the assumption, but I'm asking out of curiosity.}
    This confirms the base case. For induction, assume that the
    hypothesis is true for $A_{j-1}$. Consider processing the column $j$
    in $A$. Let $j'=_\sigma j$ be the column in $A'$ corresponding to column
    $j$ in $A$. If the algorithm {\sc ConMat} adds a column $k$ to
    column $j$, we have (i) $k <j$,  (ii) 
    % $[k]_P \leq_P [j]_P$
    $[k]_P \leq [j]_P$
    by Proposition~\ref{prop:morsepath}, and (iii) column $k$ is homogeneous. Let $k=_\sigma k'$. Then, 
    $[k']_P=[k]_P\leq_P [j]_P=[j']_P$ 
    % $[k']_P=[k]_P\leq_P [j]_P=[j']_P$ 
    and hence $k'< j'$ in $A'$.
    By inductive hypothesis,
    $A_{out}[i,k]= A'_{out}[i',k']$ for all $i=_\sigma i'$.
    Also, since $k$ is a homogeneous column,
    so is $k'$ by inductive hypothesis and their pivot rows, say $\ell$ and $\ell'$
    satisfy $\ell=_\sigma \ell'$. 
    It follows that when {\sc ConMat} processes the column $j'$, it also
    adds column $k'$ to $j'$. Therefore, {\sc ConMat} adds every column 
    to column $j'$ while processing $A'$ whose corresponding columns it adds to $j$ while processing
    $A$. Reversing the roles of $A$ and $A'$, we can claim that {\sc ConMat}
    adds every column to $j$ while processing $A$ whose corresponding columns it adds to $j'$ while
    processing $A'$. Essentially, this means that $A_{out}[i,j]=A'_{out}[i',j']$
    where $i=_\sigma i'$ and $j=_\sigma j'$. Furthermore, since the order
    of the simplices within each Morse set remains unaltered, the pivot rows of $j$ and $j'$ (if they remain homogeneous) 
    correspond to the
    same simplex completing the induction. \qed
\end{proof}

\end{document}
